% Chapitre 9 — RAG et rapports
\chapter{RAG et génération de rapports intelligents}
\label{ch:rag}

\section{Principe du RAG}

Le \textbf{Retrieval-Augmented Generation} combine :
\begin{enumerate}[leftmargin=*]
  \item \textbf{Retrieval} : recherche de passages pertinents dans une base documentaire ;
  \item \textbf{Augmentation} : injection de ces passages dans le prompt LLM ;
  \item \textbf{Generation} : rédaction du rapport ou de l'explication.
\end{enumerate}

\textbf{Règle fondamentale du projet} : le RAG enrichit le \textit{texte}, jamais les \textit{scores numériques} produits par les modèles ML.

\section{Architecture RAG locale}

\begin{itemize}[leftmargin=*]
  \item \textbf{Index} : TF-IDF (scikit-learn), matrice sparse, chunks de $\sim$850 caractères ;
  \item \textbf{Stockage} : \texttt{models/rag\_index.joblib} ;
  \item \textbf{Corpus curé} :
    \begin{itemize}
      \item \texttt{reports/knowledge\_base/ref08\_extrait.txt}
      \item \texttt{reports/knowledge\_base/regles\_metier\_microfinance.md}
      \item \texttt{reports/criteres\_score\_et\_crispdm.md}
    \end{itemize}
  \item \textbf{Exclusions} : rapports générés, fichiers demo, notes dev.
\end{itemize}

\section{Pipeline de retrieval}

\subsection{Multi-query retrieval}

Plutôt qu'une requête monolithique, le module \texttt{build\_rag\_query()} génère plusieurs sous-requêtes :
\begin{itemize}[leftmargin=*]
  \item requête REF-08 générique ;
  \item requête par niveau de risque (FAIBLE/MODERE/ELEVE) ;
  \item requête par modèle (LSTM, GraphSAGE...) ;
  \item requête par score KYC ;
  \item tokens significatifs du message utilisateur.
\end{itemize}

La fonction \texttt{retrieve\_multi()} fusionne les résultats en gardant le meilleur score par chunk.

\subsection{Seuil de pertinence}

Un score cosine minimum (\texttt{min\_score=0.06}) filtre le bruit. Si aucun chunk ne dépasse le seuil, le meilleur résultat est conservé pour ancrage REF-08.

\section{Formatage pour le LLM}

\texttt{format\_rag\_sources\_for\_prompt()} produit :
\begin{lstlisting}[caption={Exemple extrait RAG injecté au LLM}]
## Extraits documentaires (RAG)
### [1] knowledge_base/ref08_extrait.txt#0 (pertinence=0.329)
REF-08 — Spécifications Techniques...
\end{lstlisting}

Limite de contexte : 4000 caractères (configurable).

\section{Génération de rapports}

\subsection{Endpoint /report/by-cin}

Workflow :
\begin{enumerate}[leftmargin=*]
  \item Scoring classique + séquentiel + graphe (LLM off) ;
  \item Retrieval RAG multi-query ;
  \item Prompt LLM avec JSON résultats + extraits RAG ;
  \item Fallback Markdown déterministe si Ollama timeout ;
  \item Sauvegarde \texttt{reports/generated/report\_\{cin\}.md}.
\end{enumerate}

\subsection{Structure obligatoire du rapport}

\begin{enumerate}[leftmargin=*]
  \item Résumé
  \item Résultat classique
  \item Résultat séquentiel
  \item Résultat graphe
  \item Justification
  \item Recommandation
  \item Références (citations source\#chunk)
  \item Disclaimer (données synthétiques, décision humaine)
\end{enumerate}

\section{Export PDF}

Module ReportLab (\texttt{\_markdown\_to\_simple\_pdf}) :
\begin{itemize}[leftmargin=*]
  \item logo Talys en en-tête (\texttt{reports/assets/talys\_logo.png}) ;
  \item normalisation Unicode $\rightarrow$ ASCII pour compatibilité fonts ;
  \item pagination automatique.
\end{itemize}

\section{Enrichissement des explications /explain}

Les endpoints \api{/explain*} injectent désormais un contexte RAG compact dans le prompt Ollama via \texttt{\_rag\_context\_for\_scoring()}, améliorant la pertinence métier des explications courtes.

\section{Reconstruction de l'index}

\begin{lstlisting}[caption={Rebuild manuel de l'index RAG}]
python -m src.rag.index
\end{lstlisting}

Rebuild automatique si un fichier source est modifié (comparaison \texttt{mtime}).

\section{Limites et évolutions}

\begin{table}[H]
  \centering
  \caption{RAG actuel vs cible production}
  \begin{tabular}{|l|l|l|}
    \hline
    Aspect & Prototype & Production \\
    \hline
    Embedding & TF-IDF & Sentence-BERT / OpenAI \\
    \hline
    Vector DB & Joblib local & FAISS / Chroma / Pinecone \\
    \hline
    Corpus & 3 fichiers & REF-08 complet + policies \\
    \hline
    Evaluation & Manuel & RAGAS, faithfulness score \\
    \hline
  \end{tabular}
\end{table}

\section{Éthique et conformité}

\begin{itemize}[leftmargin=*]
  \item Le LLM ne doit pas inventer de chiffres absents des entrées JSON ;
  \item Les références citées sont vérifiables (source\#chunk) ;
  \item Disclaimer explicite dans chaque rapport ;
  \item Décision crédit reste humaine (comité, analyste senior).
\end{itemize}
